\PassOptionsToPackage{unicode}{hyperref}
\PassOptionsToPackage{hyphens}{url}
\documentclass[11pt,a4paper]{article}

\usepackage[T1]{fontenc}
\usepackage[utf8]{inputenc}
\usepackage[brazil]{babel}
\usepackage{lmodern}
\usepackage{microtype}

\usepackage{xcolor}
\definecolor{javaOrange}{HTML}{E76F00}
\definecolor{javaDeep}{HTML}{BF5700}
\definecolor{javaBlue}{HTML}{1565C0}
\definecolor{javaTeal}{HTML}{00695C}
\definecolor{javaAmber}{HTML}{B45309}
\definecolor{javaInk}{HTML}{2B2140}
\definecolor{javaCodeBg}{HTML}{FFF8F0}
\definecolor{javaRule}{HTML}{FFD7B5}
\definecolor{javaShade}{HTML}{FFF3E8}

\usepackage{tcolorbox}
\tcbuselibrary{skins,breakable,listings}
\usepackage{listings}
\usepackage{fvextra}

\lstdefinestyle{javastyle}{
  language=Java,
  basicstyle=\ttfamily\small,
  keywordstyle=\color{javaBlue}\bfseries,
  stringstyle=\color{javaDeep},
  commentstyle=\color{javaTeal}\itshape,
  breaklines=true,
  showstringspaces=false,
  tabsize=4,
}

\newtcblisting{javacode}{
  listing only,
  listing style=javastyle,
  breakable, enhanced,
  colback=javaCodeBg, colframe=javaDeep,
  leftrule=4pt, rightrule=0.4pt, toprule=0.4pt, bottomrule=0.4pt,
  arc=2pt, fontupper=\small,
  top=4pt, bottom=4pt, left=6pt, right=6pt,
}

\newtcbox{\inlinecode}{
  on line,
  colback=javaCodeBg, colframe=javaRule,
  boxrule=0.5pt, arc=2pt,
  fontupper=\ttfamily\small\color{javaDeep},
  boxsep=1pt, left=2pt, right=2pt, top=1pt, bottom=1pt,
}

\usepackage[margin=2.4cm, top=2.8cm, bottom=2.8cm]{geometry}

\usepackage{fancyhdr}
\pagestyle{fancy}
\fancyhf{}
\renewcommand{\headrulewidth}{0.4pt}
\renewcommand{\headrule}{\hbox to\headwidth{\color{javaRule}\leaders\hrule height \headrulewidth\hfill}}
\fancyhead[L]{\small\sffamily\color{javaDeep} Java Programming MOOC}
\fancyhead[R]{\small\sffamily\color{javaDeep} Licao 7}
\fancyfoot[C]{\small\sffamily\color{gray}\thepage}

\usepackage{titlesec}
\titleformat{\section}
  {\sffamily\Large\bfseries\color{javaOrange}}{\thesection}{0.7em}{}
  [\vspace{2pt}{\color{javaRule}\titlerule[1.2pt]}]
\titleformat{\subsection}
  {\sffamily\large\bfseries\color{javaDeep}}{\thesubsection}{0.6em}{}
\titleformat{\subsubsection}
  {\sffamily\normalsize\bfseries\color{javaTeal}}{\thesubsubsection}{0.5em}{}

\usepackage{enumitem}
\setlist[itemize]{itemsep=2pt, topsep=4pt, parsep=0pt}
\setlist[enumerate]{itemsep=2pt, topsep=4pt, parsep=0pt}

\usepackage{booktabs}
\usepackage{array}
\usepackage{colortbl}
\usepackage{longtable}

\usepackage{hyperref}
\hypersetup{
  colorlinks=true,
  linkcolor=javaBlue, urlcolor=javaBlue, citecolor=javaBlue,
  pdftitle={Programacao em Java - Licao 7: Paradigmas, Algoritmos e Exercicios Finais},
  pdfauthor={Java Programming MOOC - Universidade de Helsinque},
  pdflang={pt-BR},
}

\newtcolorbox{destaque}[1][Dica]{
  enhanced, breakable,
  colback=javaShade, colframe=javaOrange,
  leftrule=4pt, rightrule=0.4pt, toprule=0.4pt, bottomrule=0.4pt,
  arc=2pt,
  fonttitle=\sffamily\bfseries\color{javaOrange},
  title=#1,
  top=4pt, bottom=4pt,
}

\makeatletter
\newcommand{\subtitle}[1]{\gdef\@subtitle{#1}}
\renewcommand{\maketitle}{%
  \begin{center}
    {\color{javaRule}\rule{\linewidth}{1pt}}\\[6pt]
    {\Huge\sffamily\bfseries\color{javaOrange} \@title}\\[6pt]
    \ifx\@subtitle\undefined\else
      {\Large\sffamily\color{javaDeep} \@subtitle}\\[4pt]
    \fi
    {\small\sffamily\color{gray} \@author}\\[2pt]
    {\small\sffamily\color{gray} \@date}\\[6pt]
    {\color{javaRule}\rule{\linewidth}{1pt}}
  \end{center}
  \vspace{1em}
}
\makeatother

\title{Programacao em Java}
\subtitle{Licao 7: Paradigmas de Programacao, Algoritmos e Exercicios Finais}
\author{Java Programming MOOC - Universidade de Helsinque --- traducao para o portugues}
\date{2026}

\begin{document}
\maketitle
\tableofcontents
\vspace{1.5em}

\section{Objetivos de aprendizagem}

\begin{itemize}
  \item Compreender o conceito de paradigma de programacao
  \item Distinguir a programacao procedural da orientada a objetos
  \item Implementar o algoritmo de ordenacao por selecao
  \item Usar \inlinecode{Arrays.sort()} e \inlinecode{Collections.sort()}
  \item Implementar busca linear e busca binaria
  \item Distinguir metodos estaticos de metodos de instancia
  \item Aplicar os conceitos do curso em exercicios de maior porte
\end{itemize}

\section{Parte 1: Paradigmas de Programacao}

\subsection{O que e um paradigma?}

Um \textbf{paradigma de programacao} e uma maneira de pensar e estruturar a funcionalidade de um programa. Os principais sao:
\begin{enumerate}
  \item \textbf{Programacao orientada a objetos} (POO)
  \item \textbf{Programacao procedural}
  \item \textbf{Programacao funcional}
\end{enumerate}

\subsection{Comparacao: Relogio}

\textbf{Versao procedural:}

\begin{javacode}
int horas = 0, minutos = 0, segundos = 0;
while (true) {
    System.out.printf("%02d:%02d:%02d%n", horas, minutos, segundos);
    segundos++;
    if (segundos >= 60) { segundos = 0; minutos++; }
    if (minutos >= 60)  { minutos = 0;  horas++;   }
    if (horas >= 24)    { horas = 0; }
}
\end{javacode}

\textbf{Versao orientada a objetos:}

\begin{javacode}
public class Ponteiro {
    private int valor, limite;
    public Ponteiro(int limite) { this.limite = limite; this.valor = 0; }
    public void avanca() { this.valor++; if (this.valor >= this.limite) this.valor = 0; }
    public int getValor() { return this.valor; }
    public String toString() { return (this.valor < 10 ? "0" : "") + this.valor; }
}

public class Relogio {
    private Ponteiro horas, minutos, segundos;
    public Relogio() {
        this.horas = new Ponteiro(24);
        this.minutos = new Ponteiro(60);
        this.segundos = new Ponteiro(60);
    }
    public void avanca() {
        this.segundos.avanca();
        if (this.segundos.getValor() == 0) {
            this.minutos.avanca();
            if (this.minutos.getValor() == 0) this.horas.avanca();
        }
    }
    public String toString() { return horas + ":" + minutos + ":" + segundos; }
}
\end{javacode}

\begin{center}
\begin{tabular}{l l l}
\toprule
\textbf{Aspecto} & \textbf{Procedural} & \textbf{Orientado a objetos} \\
\midrule
Estado & Variaveis soltas & Encapsulado nos objetos \\
Responsabilidade & Concentrada & Distribuida \\
Reuso & Dificil & Facil \\
\bottomrule
\end{tabular}
\end{center}

\subsection{Programacao funcional em Java 8+}

\begin{javacode}
List<Integer> numeros = List.of(1, 2, 3, 4, 5, 6);
int somaDosParES = numeros.stream()
    .filter(n -> n % 2 == 0)
    .mapToInt(Integer::intValue)
    .sum();
System.out.println(somaDosParES); // 12
\end{javacode}

\section{Parte 2: Algoritmos}

\subsection{Ordenacao por selecao}

\begin{javacode}
public static void ordenacaoPorSelecao(int[] vetor) {
    for (int i = 0; i < vetor.length - 1; i++) {
        int indiceMinimo = i;
        for (int j = i + 1; j < vetor.length; j++) {
            if (vetor[j] < vetor[indiceMinimo]) indiceMinimo = j;
        }
        int temp = vetor[i];
        vetor[i] = vetor[indiceMinimo];
        vetor[indiceMinimo] = temp;
    }
}
\end{javacode}

\subsection{Ordenacao com metodos do Java}

\begin{javacode}
int[] nums = {5, 2, 8, 1, 9};
Arrays.sort(nums);
System.out.println(Arrays.toString(nums)); // [1, 2, 5, 8, 9]

ArrayList<Integer> lista = new ArrayList<>(Arrays.asList(5, 2, 8, 1));
Collections.sort(lista);
System.out.println(lista); // [1, 2, 5, 8]
\end{javacode}

\subsection{Busca linear}

\begin{javacode}
public static int buscaLinear(int[] vetor, int alvo) {
    for (int i = 0; i < vetor.length; i++) {
        if (vetor[i] == alvo) return i;
    }
    return -1;
}
\end{javacode}

\subsection{Busca binaria}

\begin{javacode}
public static int buscaBinaria(int[] vetor, int alvo) {
    int esq = 0, dir = vetor.length - 1;
    while (esq <= dir) {
        int meio = (esq + dir) / 2;
        if (vetor[meio] == alvo)      return meio;
        else if (vetor[meio] < alvo)  esq = meio + 1;
        else                          dir = meio - 1;
    }
    return -1;
}
\end{javacode}

\begin{center}
\begin{tabular}{l c c}
\toprule
\textbf{Algoritmo} & \textbf{Dados desordenados?} & \textbf{Pior caso} \\
\midrule
Busca linear  & Sim & $O(n)$ \\
Busca binaria & Nao & $O(\log n)$ \\
\bottomrule
\end{tabular}
\end{center}

\subsection{Metodos estaticos vs.\ de instancia}

\begin{javacode}
// Metodo de instancia -- acessa o estado do objeto
public void incrementar() { this.valor++; }

// Metodo estatico -- nao depende de nenhum objeto
public static int maximo(int a, int b) { return a > b ? a : b; }

// Chamada sem criar objeto:
Matematica.maximo(5, 10);
\end{javacode}

\section{Parte 3: Exercicios de Programacao Maiores}

\begin{destaque}[Dica]
Identifique os substantivos (classes) e verbos (metodos) do enunciado antes de comecar a codificar.
\end{destaque}

\subsection{Exercicio 1 --- Agenda de contatos}

Implemente uma agenda de contatos com menu interativo.

\begin{itemize}
  \item \textbf{Classe Contato:} \inlinecode{nome}, \inlinecode{telefone}, \inlinecode{email}. Getters, \inlinecode{toString()}, \inlinecode{equals()} por nome.
  \item \textbf{Classe Agenda:} \inlinecode{adicionarContato}, \inlinecode{buscarPorNome} (parcial), \inlinecode{removerContato}, \inlinecode{listarTodos} (ordem alfabetica).
  \item \textbf{main:} menu com adicionar, buscar, remover, listar, sair.
\end{itemize}

\subsection{Exercicio 2 --- Registro de notas}

\begin{itemize}
  \item \textbf{Classe Aluno:} nome, matricula, notas (0--10); \inlinecode{calcularMedia()}, \inlinecode{melhorNota()}, \inlinecode{piorNota()}, \inlinecode{aprovado(double)}.
  \item \textbf{Classe Turma:} \inlinecode{mediaGeral()}, \inlinecode{melhorAluno()}, \inlinecode{listarAprovados(double)}.
  \item \textbf{main:} crie turma com 5 alunos, 3 notas cada; exiba media geral, melhor aluno, aprovados ($\geq 6$).
\end{itemize}

\subsection{Exercicio 3 --- Loja simples}

\begin{itemize}
  \item \textbf{Classe Produto:} nome, preco, estoque; \inlinecode{vender(int)}, \inlinecode{repor(int)}.
  \item \textbf{Classe Carrinho:} \inlinecode{HashMap<Produto, Integer>}; \inlinecode{adicionar}, \inlinecode{remover}, \inlinecode{calcularTotal()}.
  \item \textbf{main:} simule cliente adicionando 3 itens e finalizando a compra.
\end{itemize}

\section{Conclusao do Java Programming I}

\begin{center}
\begin{longtable}{c l}
\toprule
\textbf{Parte} & \textbf{Topico} \\
\midrule
1 & Primeiros passos: variaveis, tipos, condicionais \\
2 & Repeticao com \texttt{while} e \texttt{for}, metodos \\
3 & Listas (\texttt{ArrayList}), arrays, Strings \\
4 & Orientacao a objetos: classes, objetos, encapsulamento \\
5 & OOP avancada: sobrecarga, referencias, composicao \\
6 & Listas de objetos, separacao UI/logica, testes unitarios \\
7 & Paradigmas, algoritmos de ordenacao e busca \\
\bottomrule
\end{longtable}
\end{center}

No \textbf{Java Programming II}, voce aprendera: heranca, polimorfismo, interfaces, generics, streams, lambdas, excecoes avancadas e muito mais.

\bigskip
\begin{center}
  {\large\sffamily\bfseries\color{javaOrange} Parabens e continue programando!}
\end{center}

\end{document}
